\section{Perfilado de la aplicación}

\subsection{Objetivo del perfilado}

El perfilado responde una pregunta distinta a la del benchmark. El benchmark mide \emph{cuánto}; el perfilado busca \textbf{explicar por qué}. Concretamente, en este trabajo se busca:

\begin{enumerate}
    \item Identificar \textbf{dónde se va el tiempo de CPU} de la aplicación (perfilado de CPU).
    \item Identificar \textbf{cuánta presión sobre la memoria} genera el kernel y si el acceso a la caché es eficiente (perfilado de memoria).
    \item Explicar el \textbf{escalamiento} observado: por qué el speedup con OpenMP es alto pero no se acerca al ideal, y por qué la transposición de $B$ produce una ganancia mayor que la paralelización.
\end{enumerate}

\subsection{Perfilado de CPU con GProf}

GProf es un perfilador por \textbf{muestreo} que se integra con el compilador. El procedimiento, tal como lo presentó el profesor, tiene tres pasos:

\begin{enumerate}
    \item \textbf{Compilar con instrumentación}: la bandera \texttt{-pg} le indica al compilador que inserte llamadas a funciones de conteo en el programa.
    \item \textbf{Ejecutar}: el programa instrumentado corre normalmente y al terminar escribe un archivo \texttt{gmon.out} con los contadores recolectados.
    \item \textbf{Generar el reporte}: \texttt{gprof binario gmon.out} convierte ese archivo binario en un reporte legible en texto plano.
\end{enumerate}

El reporte contiene el \emph{flat profile}, donde cada función aparece con el porcentaje de tiempo propio, los segundos acumulados, el número de llamadas y el tiempo medio por llamada.

Se aplicó a las tres variantes del kernel: secuencial, transpuesta y con OpenMP. Hay una \textbf{limitación importante con OpenMP}: GProf solo instrumenta el hilo principal, de modo que el tiempo de la región paralela se contabiliza de forma engañosa (aparece atribuido a la función que la contiene y no a la función de trabajo real). Por eso el reporte de GProf es concluyente para las variantes secuenciales y se reporta con la advertencia explícita para la variante con OpenMP.

\subsection{Perfilado de memoria con Cachegrind}

Para el eje de memoria se usó \textbf{Cachegrind}, la herramienta de simulación de caché que forma parte de Valgrind. Cachegrind reproduce el comportamiento de una jerarquía de caché durante la ejecución y reporta, para el programa completo y para cada función:

\begin{itemize}
    \item \textbf{Instrucciones ejecutadas} (\textit{I refs}).
    \item \textbf{Referencias a datos} (\textit{D refs}), separadas en lecturas y escrituras.
    \item \textbf{Fallos de caché} (\textit{misses}) en L1 de datos (\textit{D1}) y en el último nivel (\textit{LL}).
    \item \textbf{Tasa de fallo} (\textit{miss rate}) de lectura y de escritura.
\end{itemize}

Para activar la simulación de caché hay que pasar la opción \texttt{--cache-sim=yes}; sin ella Cachegrind solo reporta el conteo de instrucciones.

Valgrind no venía instalado en el sistema, pero no hizo falta permiso de administrador: el script \texttt{experiments/bootstrap\_tools.sh} descarga el paquete con \texttt{apt-get download} (que no requiere \texttt{sudo}) y lo extrae con \texttt{dpkg-deb -x} en una carpeta local del proyecto, ajustando la variable \texttt{VALGRIND\_LIB} para que encuentre sus archivos de soporte.

\subsection{Footprint de memoria con \texttt{getrusage}}

Para completar el eje de memoria con la medida del \emph{tamaño} del espacio de memoria utilizado (y no solo su eficiencia), se midió el \textbf{max RSS} (\textit{maximum resident set size}), que es la cantidad máxima de memoria física que el proceso tuvo residente simultáneamente. Se obtuvo con la llamada al sistema \texttt{getrusage(RUSAGE\_CHILDREN, ...)}, envuelta en un pequeño programa auxiliar (\texttt{experiments/maxrss.c}) que ejecuta la aplicación y reporta el pico de memoria del hijo. Esta medida no depende de ninguna herramienta externa.

\subsection{Alcance del perfilado}

El perfilado \textbf{caracteriza la aplicación}: dice dónde se va el tiempo y cuánta presión sobre la memoria genera el kernel. El perfilado de CPU (GProf) se corre en las dos máquinas, mientras que el detalle fino de caché (Cachegrind) se midió en la máquina A, que es donde se corrió el experimento de desempeño principal.

La comparación entre las dos máquinas se hace por otra vía: con el \textbf{benchmark}, que sí se ejecutó en los dos equipos (sección de comparación entre máquinas).

\subsection{Tamaño de N usado para el perfilado}

Cachegrind simula la caché instrucción por instrucción, lo que lo hace entre 50 y 100 veces más lento que la ejecución nativa. Por eso el perfilado se hizo con $N = 500$ en lugar de $N = 4000$. Esta limitación es aceptable: la \emph{tasa} de fallo de caché es una propiedad del \emph{patrón de acceso}, y el patrón de acceso del kernel es el mismo para cualquier $N$ (solo cambia la cantidad de trabajo). Es más, se verificó que la tasa de fallo se mantiene estable entre $N = 500$ y $N = 1000$.
